This analysis has charted a complete narrative arc in humanity's understanding of the a dream state.
The journey began with humans as \textit{passive recipients} of "divine" messages (Ramanujan, Sec 1).
It progressed, through scientific inquiry, to a model of the brain as a \textit{passive trainee} in an automatic, "offline" computational process (the "AI training," Sec 3).
It now concludes with the human as an \textit{active agent*} in their own cognitive development.

The Dormio (TDI) technology represents the first step, allowing us to deliberately \textit{induce} the creative synthesis that people like Kekulé experienced by chance \citep{haarhorowitz2020dormio}. But lucid dreaming represents the "final frontier." If sleep is the "offline training phase" where the brain consolidates its past and simulates its future, lucid dreaming is the act of gaining conscious, administrative control over that phase. The dreamer can "rehearse" motor skills \citep{dresler2011motor} or "resolve" psychological conflicts \citep{voss2009lucid}, actively curating their own mental "training data." This completes the journey: from believing dreams are messages \textit{from} a god, to understanding they are a training simulation \textit{by} the brain, to finally seizing the controls and \textit{piloting} that simulation ourselves.